\section{Advantage Actor-Critic (A2C)}
    \subsection{Main Idea}
        The core idea of A2C is to reduce the variance by introducing the baseline.
        The following equation enable us to use the baseline when estimating state values.
        \begin{equation}
            \E_{S \sim \eta, A \sim \pi}[\nabla_\theta \log{\pi_\theta(A|S)}q_\pi(S,A)] = \E_{S \sim \eta, A \sim \pi}[\nabla_\theta\log{\pi_\theta(A|S)}(q_\pi(S,A)-b(S))]
        \end{equation}
        where $b(S)$ is a scalar function of $S$.
        \begin{proof}
            \begin{align}
                \E_{S \sim \eta, A \sim \pi}\Big[ \nabla_\theta \log{\pi_\theta\log{\pi(A|S)}b(S)} \Big]
                    &= \sum_{s \in \mathcal{S}} \eta(s) \sum_{a \in \mathcal{A}(s)}\pi_\theta(a|s)\nabla_\theta \log{\pi_\theta(a|s)}b(s) \\
                    &= \sum_{s \in \mathcal{S}} \eta(s) \sum_{a \in \mathcal{A}(s)}\nabla_\theta \pi_\theta(a|s) b(s) \\
                    &= \sum_{s \in \mathcal{S}} \eta(s)b(s) \sum_{a \in \mathcal{A}(s)} \nabla_\theta \pi_\theta(a|s) \\
                    &= \sum_{s \in \mathcal{S}} \eta(s)b(s) \nabla_\theta \sum_{a \in \mathcal{A}(s)}\pi_\theta(a|s) \\
                    &= \sum_{s \in \mathcal{S}} \eta(s)b(s) \nabla_\theta 1 = 0.
            \end{align}
        \end{proof}

        The baseline reduces the variance of estimation of gradients. In particular, let
        \begin{equation}
            X(S,A) = \nabla_\theta \log{\pi_\theta(a|s)}[q_\pi(S,A)-b(S)]
        \end{equation}
        Then, the true gradient is $\E[X(S,A)]$ and we estimate it with samples $x \sim X$. To better estimate the gradient,
        it would be favorable to make the variance of $X$ small. The theoretically optimal baseline that minimizes $\text{var}(X)$ is
        \begin{equation}
            b^*(s) = \frac{\E_{A \sim \pi}\big[ \|\nabla_\theta \log_\theta(A|s)\|^2 q_\pi(s,A) \big]}
            {\E_{A \sim \pi}\big[ \|\nabla_\theta\log{\pi(A|s)}\|^2 \big]}
        \end{equation}
        proof can be found at \cite{Zhao2025}.
        But it is too complex to be useful in practice. In practice, we use the suboptimal version without the logarithm term, which is actually a value function, intersetingly.
        \begin{equation}
            b(s) = \E_{A \sim \pi}[q_\pi(s,A)] = v_\pi(s)
        \end{equation}

    \subsection{Pseudocode}
        \Algorithm{a2c}
        {Advantage Actor-Critic Algorithm}
        {
            \inputitem{$\pi_\theta$}{a policy function with initial parameter $\theta_0$} \\
            \inputitem{$v_w(s)$}{a value function with initial parameter $w_0$}
        }
        {
            \inputitem{$\pi^*$}{(sub)optimal policy}
        }
        {
            \ForEach{time step $t$} {
                generate and take an action $a_t$ following $\pi_\theta$\;
                observe $r_{t+1}$ and $s_{t+1}$.\;
                \tcp{advantage}
                $\delta_t \leftarrow r_{t+1} + \gamma v_{w_t}(s_{t+1}) - v_{w_t}(s_t)$\;
                \tcp{actor (policy update)}
                $\theta_{t + 1} \leftarrow \theta_t + \alpha_\theta \delta_t \nabla_\theta \log{\pi_{\theta_t}(a_t|s_t)}$\;
                \tcp{critic (value update)}
                $w_{t+1} \leftarrow w_t + \alpha_w \delta_t \nabla_w v_{w_t}(s_t)$\; 
            }
            \Return{$\pi$}
        }
        
        Using the baseline $b(s) = v_\pi(s)$, we get the following update equation:
        \begin{equation}
            \theta_{t+1} = \theta_t + \alpha \nabla_\theta \log{\pi_{\theta_t}(a_t|s_t)}[q_t(s_t,a_t) - v(s_t)]
        \end{equation}

        If $q_t(s_t,a_t)$ and $v_t(s_t)$ were estimated by Monte Carlo, the algorithm is called \textbf{REINFORCE with a baseline}.
        If estimated by TD learining, the algorithm~\ref{alg:a2c} is called \textbf{A2C} or \textbf{TD actor-critic}.
        Note that algorithm~\ref{alg:a2c} estimates $q_t(s_t,a_t) - v_t(s_t)$ with TD error, which is $r_{t+1} + \gamma v_t(s_{t+1} - v_t(s_t))$.
    
    \subsection{Implementation Details}
        We can use n-step update to balance between bias and variance of estimation.
        \begin{equation}
            G_t^{(n)} = r_{t + 1} + \gamma r_{t+2} + \gamma^2 r_{t+3} + \dots + \gamma^n v_t(s_{n+1})
        \end{equation}

        To stop the policy from being almost deterministic too quickly, we can add \textbf{entropy bonus} at the loss function (\cite{pmlr-v48-mniha16}, \cite{WILLIAMS1991}).
        Specifically, the loss function is
        \begin{equation}
            \mathcal{L} = \frac{1}{N}\sum_{i = 0}^{N-1} \Big[-\log{\pi_\theta(a_{t+i}|s_{t+i})}A_{t+i} + c_v(v_\phi(s_{t+i}) - G_{t+i})^2 - c_e \mathcal{H}(\pi_\theta(\cdot|s_{t+i}))\Big]
        \end{equation}
        where 
        \begin{align}
            \mathcal{H}(\pi(\cdot|s)) &= -\int \pi(a|s)\log{\pi(a|s)}da \\
            \mathcal{H} &= -\sum_a \pi(a|s)\log{\pi(a|s)},
        \end{align}
        $A$ and $G$ are advantage and return estimate, respectively.

        The actor and critic can share parameters, when they are using the encoder. The can use different encoder, each of their own, or share one encoder.
    \subsection{Pros and Cons}
    Pros:
    \begin{itemize}
        \item Requires only one networks, opposed to QAC;
        \item Variance is decreased, opposed to the original REINFORCE without baseline;
        \item Appliable to both discrete and continuous actions;
    \end{itemize}
    Cons:
    \begin{itemize}
        \item If the baseline is incorrect, it can be biased;
        \item Sample efficiency is low, since it is an on-policy method;
        \item Vulnerable to big, sudden change of parameters. \textbf{PPO} was designed to resolve this;
    \end{itemize}